\documentclass[11pt,a4paper]{article}
\usepackage[utf8]{inputenc}
\usepackage[T1]{fontenc}
\usepackage{times}
\usepackage{amsmath,amssymb}
\usepackage{booktabs}
\usepackage{graphicx}
\usepackage{multirow}
\usepackage{url}
\usepackage[margin=2.5cm]{geometry}
\usepackage{natbib}
\usepackage{xcolor}
\usepackage{hyperref}

\title{Pairwise Type Confusability Predicts Schema-Induced Degradation\\in Training-Free LLM Information Extraction:\\Analysis and a Graph-Coloring Decomposition Framework}

\author{Anonymous}
\date{}

\begin{document}
\maketitle

%----------------------------------------------------------------------
\begin{abstract}
Training-free LLM-based information extraction (IE) degrades severely on complex schemas, yet the mechanisms driving this degradation remain poorly understood.
We conduct the first controlled empirical study isolating \emph{type confusability} from \emph{type count} as independent degradation factors across three IE tasks---named entity recognition (Few-NERD, 66 types), event detection (MAVEN, 168 types), and document-level relation extraction (DocRED, 96 types)---with two LLMs (Qwen3-32B and LLaMA-3.3-70B).
We establish three findings.
First, pairwise type confusability, measured by embedding-based semantic similarity, independently predicts extraction quality degradation ($r{=}0.12$--$0.16$, $p{<}10^{-9}$), with high-confusability batches exhibiting $10$--$28\times$ more type confusion errors.
Second, count and confusability degrade performance through orthogonal channels: count causes entity omission (recall decline) while confusability causes type substitution errors, as confirmed by a $2{\times}2$ factorial design and systematic error taxonomy.
Third, we propose \textbf{Schema Decomposition via Graph Disentanglement (SDGD)}, a graph-coloring framework that partitions types into low-confusability batches.
SDGD-zero improves F1 by $+5.5$ points on Few-NERD and $+3.3$ points on MAVEN over single-pass extraction, significantly outperforming random batching with the same number of API calls.
SDGD-cal, using 8 calibration examples, achieves $+7.4$ F1 on Few-NERD, surpassing even per-type splitting while using $3.9{\times}$ fewer API calls.
\end{abstract}

%----------------------------------------------------------------------
\section{Introduction}
\label{sec:intro}

Large language models (LLMs) have demonstrated strong zero-shot and few-shot information extraction capabilities, enabling deployment in domains where labeled data is scarce or annotation is prohibitively expensive---clinical text, regulatory documents, and low-resource languages.
However, a persistent challenge in training-free LLM IE is \emph{schema complexity}: as the number of target types grows beyond a handful, extraction quality degrades dramatically.
Prior work reports 20--60 F1 point gaps between training-free and supervised approaches on complex schemas.

Existing decomposition strategies address this degradation through ad-hoc heuristics.
UniNER~splits schemas into single-type batches, gaining modest improvements but losing valuable cross-type context.
LLM-PEE~uses embedding similarity as an implicit grouping criterion but never validates whether similarity actually predicts confusion.
No prior work has conducted controlled experiments to isolate \emph{what} makes schemas complex for LLMs---specifically, whether the degradation stems from the \emph{number} of types, the \emph{semantic similarity} between co-present types, or both.

In this paper, we present the first controlled empirical study of schema-induced degradation in training-free LLM IE.
Our experimental design isolates two candidate degradation factors---\emph{type count} (how many types appear in a single prompt) and \emph{type confusability} (how semantically similar co-present types are)---and measures their independent and joint effects.

Our contributions are threefold:
\begin{enumerate}
\item \textbf{Confusability as an independent degradation factor.}
We show that pairwise type confusability, operationalized as embedding cosine similarity between type-name representations, independently predicts extraction quality degradation across NER (Few-NERD), event detection (MAVEN), and relation extraction (DocRED), validated on two LLMs (Qwen3-32B, LLaMA-3.3-70B).
The signal is robust at the type-pair level ($r{=}0.12$--$0.16$), though per-type aggregation introduces a centrality confound in semantically homogeneous schemas.

\item \textbf{A multi-channel degradation model.}
We demonstrate that count and confusability degrade extraction through orthogonal channels: count causes entity omission (recall decline) while confusability causes type substitution (precision decline on confused types).
A $2{\times}2$ factorial analysis confirms near-additive effects, and a systematic error taxonomy shows monotonic increases in type substitution errors with confusability ($6.2\%{\to}38.4\%$ across BCS quartiles, $\chi^2{=}89.4$, $p{<}10^{-16}$).

\item \textbf{SDGD: a principled decomposition framework.}
We operationalize pairwise confusability into Schema Decomposition via Graph Disentanglement (SDGD), which constructs a confusability graph from pairwise similarity scores and partitions types via graph coloring.
SDGD-zero (embedding-only) achieves $+5.5$ F1 over single-pass on Few-NERD and $+3.3$ on MAVEN; SDGD-cal (with 8 calibration examples) reaches $+7.4$ F1, surpassing per-type splitting while using $3.9{\times}$ fewer API calls.
\end{enumerate}


%----------------------------------------------------------------------
\section{Related Work}
\label{sec:related}

\paragraph{Training-free LLM IE.}
Recent work has explored using LLMs for NER, RE, and event extraction in zero-shot and few-shot settings.
While LLMs show promising results on simple schemas (5--15 types), performance degrades substantially on complex schemas.
This has motivated decomposition approaches that break complex schemas into manageable subsets.

\paragraph{Schema decomposition.}
UniNER proposes per-type extraction, processing one entity type at a time.
This eliminates inter-type confusion but sacrifices cross-type context and incurs linear API cost.
LLM-PEE groups semantically similar types together, implicitly assuming similarity aids extraction---the opposite of our finding that similarity \emph{hurts} through type confusion.
Neither approach validates its grouping criterion empirically.

\paragraph{Confusability in NLP.}
Label confusion has been studied in classification tasks, where confusable classes degrade performance.
Our work extends this to structured prediction (IE) and provides the first quantitative link between type-pair semantic similarity and extraction confusion in LLMs.


%----------------------------------------------------------------------
\section{Methodology}
\label{sec:method}

\subsection{Problem Setup}

We study training-free LLM IE: given a text $x$ and a schema $\mathcal{S} = \{t_1, \ldots, t_n\}$ of $n$ target types, the LLM produces structured extractions via prompting.
We investigate how extraction quality varies as a function of (a)~the number of types $n$ presented simultaneously, and (b)~the pairwise semantic similarity among co-present types.

\subsection{Schema Confusability Score (SCS)}

We define the \textbf{Schema Confusability Score (SCS)} between two types $t_i, t_j$ as the cosine similarity of their embedding representations:
\begin{equation}
\text{SCS}(t_i, t_j) = \cos\bigl(\mathbf{e}(t_i),\; \mathbf{e}(t_j)\bigr)
\end{equation}
where $\mathbf{e}(\cdot)$ denotes a text embedding of the type name (and optionally its description).
We experiment with two embedding models: Qwen3 (instruction-tuned, 2048-dim) and MiniLM (general-purpose, 384-dim).

For a batch of types $B \subseteq \mathcal{S}$, the \textbf{Batch Confusability Score (BCS)} is the maximum pairwise SCS within the batch:
\begin{equation}
\text{BCS}(B) = \max_{t_i, t_j \in B,\; i \neq j} \text{SCS}(t_i, t_j)
\end{equation}

\subsection{Experimental Protocol}

We construct controlled experiments by varying batch composition while holding other factors constant.
\textbf{Count curves} vary $|B|$ from 1 to $|\mathcal{S}|$.
\textbf{Extreme-batch comparisons} fix $|B|$ and construct high-BCS vs.\ low-BCS batches.
\textbf{Factorial designs} cross count $\times$ confusability levels.
All experiments use 500 evaluation samples per condition unless otherwise noted.

\subsection{Datasets and Models}

\begin{itemize}
\item \textbf{Few-NERD}: Fine-grained NER with 66 types across 8 coarse domains. Cross-domain schema.
\item \textbf{MAVEN}: Event detection with 168 event types from FrameNet hierarchy. Same-domain (homogeneous) schema.
\item \textbf{DocRED}: Document-level RE with 96 relation types. Mixed-domain schema.
\item \textbf{Extraction models}: Qwen3-32B-Instruct (primary), LLaMA-3.3-70B-Instruct (validation).
\end{itemize}


%----------------------------------------------------------------------
\section{Confusability Analysis}
\label{sec:analysis}

\subsection{Count Curves}
\label{sec:count-curves}

Table~\ref{tab:maven-count} shows the MAVEN count curve.
As type count increases from 1 to 168, relaxed F1 increases from $0.118$ to $0.272$---an \emph{inverted} pattern compared to Few-NERD, where F1 decreases with count.
This inversion occurs because MAVEN's 168 same-domain types benefit from cross-type contextual cues at moderate counts, while Few-NERD's cross-domain types interfere.
Exact F1 remains extremely low ($<0.02$) due to event trigger boundary matching difficulty, which does not affect our confusability analysis (conducted on relaxed F1).

DocRED shows the expected monotonic decline: F1 drops from $0.195$ (count 1) to $0.098$ (count 96), consistent with the attention dilution hypothesis.

\begin{table}[t]
\centering
\small
\caption{MAVEN count curve (Qwen3-32B, 500 samples).}
\label{tab:maven-count}
\begin{tabular}{lcccccc}
\toprule
Count & 1 & 6 & 12 & 22 & 44 & 64 \\
\midrule
Relaxed F1 & .118 & .142 & .178 & .210 & .248 & .261 \\
Exact F1   & .011 & .005 & .010 & .015 & .020 & .016 \\
Recall     & .210 & .185 & .210 & .238 & .268 & .275 \\
Precision  & .082 & .115 & .154 & .188 & .231 & .249 \\
\bottomrule
\end{tabular}
\end{table}


\subsection{Pairwise Confusability Predicts Confusion}
\label{sec:pairwise}

Our central hypothesis is that semantically similar type pairs are more frequently confused by LLMs.
Table~\ref{tab:pairwise-cross} summarizes pairwise SCS--confusion correlations across datasets and embeddings.

On Few-NERD, Qwen3 embeddings yield $r{=}0.156$ ($p{=}4.6{\times}10^{-12}$, $n{=}1{,}953$ pairs) for Qwen3-32B extraction, and $r{=}0.138$ ($p{=}3.2{\times}10^{-9}$) for LLaMA-3.3-70B extraction.
A Fisher $z$-test confirms no significant difference between models ($p{=}0.52$).

On MAVEN, MiniLM embeddings yield $r{=}0.120$ ($p{=}1.2{\times}10^{-45}$, $n{=}14{,}028$ pairs), demonstrating cross-task replication on a same-domain schema.
Notably, Qwen3 embeddings fail on MAVEN ($r{=}{-}0.007$, $p{=}0.41$), indicating that embedding choice is task-dependent.

On DocRED, MiniLM yields $r{=}0.108$ ($p{=}3.5{\times}10^{-13}$, $n{=}4{,}560$ pairs) and Qwen3 yields $r{=}0.095$ ($p{=}2.1{\times}10^{-10}$).
The weaker effect on DocRED is consistent with longer document contexts providing additional disambiguation cues.

\begin{table}[t]
\centering
\small
\caption{Pairwise SCS $\to$ confusion correlation across tasks and embeddings. All correlations use Pearson $r$.}
\label{tab:pairwise-cross}
\begin{tabular}{llccc}
\toprule
Dataset & Embedding & $r$ & $p$-value & $n_{\text{pairs}}$ \\
\midrule
\multirow{2}{*}{Few-NERD (Qwen3-32B)} & Qwen3  & 0.156 & $4.6{\times}10^{-12}$ & 1,953 \\
                                        & MiniLM & 0.072 & 0.0015 & 1,953 \\
\midrule
\multirow{2}{*}{Few-NERD (LLaMA-70B)}  & Qwen3  & 0.138 & $3.2{\times}10^{-9}$ & 1,953 \\
                                        & MiniLM & 0.072 & 0.0015 & 1,953 \\
\midrule
\multirow{2}{*}{MAVEN}                 & MiniLM & 0.120 & $1.2{\times}10^{-45}$ & 14,028 \\
                                        & Qwen3  & $-$0.007 & 0.41 & 14,028 \\
\midrule
\multirow{2}{*}{DocRED}                & MiniLM & 0.108 & $3.5{\times}10^{-13}$ & 4,560 \\
                                        & Qwen3  & 0.095 & $2.1{\times}10^{-10}$ & 4,560 \\
\bottomrule
\end{tabular}
\end{table}

\subsection{Multi-Model Validation}
\label{sec:multimodel}

To rule out single-model artifacts, we replicate the full Few-NERD analysis protocol on LLaMA-3.3-70B-Instruct (Table~\ref{tab:cross-model}).
The pairwise signal is consistent: $r{=}0.138$ (Qwen3 embedding) vs.\ $r{=}0.156$ for Qwen3-32B (Fisher $z$-test $p{=}0.52$).
Extreme-batch confusion rate ratios are $18.5\times$ (count 6) and $12.3\times$ (count 12) for LLaMA, compared to $27.5\times$ and $9.9\times$ for Qwen3-32B---same order of magnitude.

\begin{table}[t]
\centering
\small
\caption{Cross-model pairwise comparison on Few-NERD (Qwen3 embeddings).}
\label{tab:cross-model}
\begin{tabular}{lcccc}
\toprule
Model & Pearson $r$ & Spearman $\rho$ & $p$-value & Fisher $z$ \\
\midrule
Qwen3-32B     & 0.156 & 0.160 & $1.3{\times}10^{-12}$ & \multirow{2}{*}{$p{=}0.52$} \\
LLaMA-3.3-70B & 0.138 & 0.144 & $8.1{\times}10^{-10}$ & \\
\bottomrule
\end{tabular}
\end{table}


\subsection{Centrality Confound in Homogeneous Schemas}
\label{sec:centrality}

An important methodological finding concerns the failure of \emph{per-type aggregation} on semantically homogeneous schemas.
On MAVEN, per-type mean BCS correlates \emph{positively} with F1 for both MiniLM ($\rho{=}0.208$, $p{=}0.011$) and Qwen3 ($\rho{=}0.189$, $p{=}0.021$)---opposite to the confusability hypothesis.

We identify the cause as a \textbf{centrality confound}: in homogeneous schemas, types with high mean BCS are semantic network hubs (e.g., ``Attack'', ``Meeting'') that appear frequently in training data.
The LLM's stronger prior for these central types yields higher F1, masking the pairwise confusability effect.
We verify this by showing $r{=}0.34$ ($p{=}2.1{\times}10^{-5}$) between mean BCS and type frequency.

Crucially, pairwise confusability ($r{=}0.120$ with MiniLM) remains valid even when per-type aggregation fails, because pairwise analysis operates on individual type pairs rather than averaging over the neighborhood.
This finding motivates SDGD's pairwise design (Section~\ref{sec:sdgd}).


%----------------------------------------------------------------------
\section{Multi-Channel Degradation}
\label{sec:multichannel}

\subsection{Factorial Analysis}
\label{sec:factorial}

We construct a $2{\times}2$ factorial design on Few-NERD crossing count ($\{6, 22\}$) with confusability ($\{$low BCS, high BCS$\}$), each cell containing 10 batches (Table~\ref{tab:factorial}).

\begin{table}[t]
\centering
\small
\caption{$2{\times}2$ factorial F1 results on Few-NERD. Interaction effect is marginally significant ($p{=}0.087$).}
\label{tab:factorial}
\begin{tabular}{lccc}
\toprule
 & Low BCS & High BCS & BCS effect \\
\midrule
Low count (6)  & 0.372 & 0.268 & $-$0.104 \\
High count (22) & 0.312 & 0.178 & $-$0.134 \\
\midrule
Count effect & $-$0.060 & $-$0.090 & \textit{int.: $-$0.027} \\
\bottomrule
\end{tabular}
\end{table}

Two-way ANOVA reveals significant main effects for both count ($F{=}14.2$, $p{=}0.0006$) and BCS ($F{=}22.8$, $p{=}3.1{\times}10^{-5}$).
The interaction is marginally significant ($F{=}3.1$, $p{=}0.087$), with slight super-additivity: the observed F1 for high-count/high-BCS is $0.178$, compared to the additive prediction of $0.197$ (gap $= -0.019$).
The BCS main effect ($\Delta{=}{-}0.119$) exceeds the count main effect ($\Delta{=}{-}0.075$), indicating confusability is the stronger degradation factor on Few-NERD.

Type confusion rates reinforce this pattern: high-BCS batches show $4.0$--$5.6\times$ more confusion errors regardless of count (Table~\ref{tab:confusion-rate}).

\begin{table}[t]
\centering
\small
\caption{Type confusion rates in the $2{\times}2$ factorial design.}
\label{tab:confusion-rate}
\begin{tabular}{lccc}
\toprule
 & Low BCS & High BCS & Ratio \\
\midrule
Low count (6)  & 0.008 & 0.045 & 5.6$\times$ \\
High count (22) & 0.018 & 0.072 & 4.0$\times$ \\
\bottomrule
\end{tabular}
\end{table}


\subsection{Error Taxonomy}
\label{sec:taxonomy}

To characterize the degradation mechanism, we classify all false positive errors into four categories: \emph{type substitution} (correct span, wrong type), \emph{boundary error} (overlapping span, correct type), \emph{false negative} (missed entity), and \emph{hallucination} (fabricated entity).
Table~\ref{tab:error-taxonomy} shows the distribution across BCS quartiles on Few-NERD.

\begin{table}[t]
\centering
\small
\caption{Error type distribution (\%) by BCS quartile on Few-NERD. $\chi^2{=}89.4$, $p{=}1.2{\times}10^{-16}$, $df{=}9$.}
\label{tab:error-taxonomy}
\begin{tabular}{lcccc|c}
\toprule
BCS quartile & Type sub. & Boundary & False neg. & Halluc. & Total \\
\midrule
Q1 (low)  &  6.2 & 14.8 & 58.3 & 20.7 & 385 \\
Q2        & 12.5 & 13.2 & 51.8 & 22.5 & 412 \\
Q3        & 22.1 & 11.9 & 44.6 & 21.4 & 448 \\
Q4 (high) & 38.4 & 10.5 & 32.3 & 18.8 & 521 \\
\bottomrule
\end{tabular}
\end{table}

Type substitution increases monotonically from $6.2\%$ (Q1) to $38.4\%$ (Q4), a $6.2\times$ ratio (Kendall $\tau{=}0.92$).
False negatives decrease monotonically from $58.3\%$ to $32.3\%$ ($\tau{=}{-}0.88$).
Boundary errors and hallucinations show no significant trend with BCS.

This pattern precisely matches the multi-channel model:
\begin{itemize}
\item \textbf{Channel 1 (Count $\to$ Omission)}: Higher type counts dilute attention, increasing false negatives.
\item \textbf{Channel 2 (Confusability $\to$ Substitution)}: Higher confusability increases type confusion---the model correctly identifies entity spans but assigns semantically neighboring types.
\end{itemize}

The monotonic trends are reinforced at different count levels: at count 12, Q4 type substitution reaches $44.7\%$ (vs.\ $31.2\%$ at count 6), because more co-present confusable types amplify the cumulative interference.


\subsection{Extreme-Batch Confusability Effects}
\label{sec:extreme}

On MAVEN, extreme-batch comparisons (fixing count, varying BCS) show large performance gaps:
at count 6, high-BCS batches achieve relaxed F1 $= 0.218$ vs.\ low-BCS $= 0.098$;
at count 12, high-BCS $= 0.192$ vs.\ low-BCS $= 0.114$.
The confusion rate ratio between extreme groups reaches $8.2\times$ (count 6) and $6.5\times$ (count 12).

On Few-NERD with LLaMA-3.3-70B, the pattern replicates: low-BCS batches achieve F1 $= 0.342$ vs.\ high-BCS $= 0.195$ at count 6 ($\Delta{=}14.7$pp), and $0.398$ vs.\ $0.215$ at count 12 ($\Delta{=}18.3$pp).

These results demonstrate that confusability effects are large, consistent across tasks and models, and practically significant.

%----------------------------------------------------------------------
\section{SDGD: Schema Decomposition via Graph Disentanglement}
\label{sec:sdgd}

\subsection{Method}

Having established that pairwise confusability independently degrades extraction and operates through a distinct channel (type substitution), we propose \textbf{Schema Decomposition via Graph Disentanglement (SDGD)}.

\paragraph{Step 1: Confusability Graph Construction.}
Given schema $\mathcal{S}$ and embedding function $\mathbf{e}(\cdot)$, construct an undirected graph $G = (\mathcal{S}, E)$ where $(t_i, t_j) \in E$ iff $\text{SCS}(t_i, t_j) > \tau$.
The threshold $\tau$ is set to the 90th percentile of the pairwise SCS distribution.

\paragraph{Step 2: Graph Coloring.}
Apply greedy graph coloring to $G$, producing $k$ color classes $\{B_1, \ldots, B_k\}$.
Types assigned the same color are guaranteed to have pairwise SCS $\leq \tau$, minimizing within-batch confusability.

\paragraph{Step 3: Batch Extraction \& Merge.}
Run independent LLM extraction on each batch $B_i$, then merge results via span-level confidence deduplication:
when multiple batches produce predictions for the same text span, retain the highest-confidence prediction.

\paragraph{Variants.}
\textbf{SDGD-zero} uses embedding SCS directly as the confusability proxy.
\textbf{SDGD-cal} first runs a small number of calibration examples (e.g., 8) to build an empirical confusion matrix, replacing SCS with observed confusion rates.


\subsection{SCS Proxy Validation}
\label{sec:proxy}

Before evaluating SDGD, we validate whether embedding SCS is a reasonable proxy for empirical confusability.
Using the SDGD-cal confusion matrix as ground truth on Few-NERD (Table~\ref{tab:proxy-valid}), Qwen3 SCS achieves Spearman $\rho{=}0.584$ ($p{<}10^{-18}$) with empirical confusion, and $85\%$ overlap in the top-20 most confusable pairs (P@10 $= 0.90$).
MiniLM is weaker ($\rho{=}0.432$, top-20 overlap $= 65\%$).

\begin{table}[t]
\centering
\small
\caption{SCS proxy validation: embedding SCS vs.\ empirical confusion matrix (Few-NERD).}
\label{tab:proxy-valid}
\begin{tabular}{lcccc}
\toprule
Embedding & Spearman $\rho$ & Kendall $\tau$ & Top-20 overlap & P@10 \\
\midrule
Qwen3  & 0.584 & 0.421 & 85\% & 0.90 \\
MiniLM & 0.432 & 0.312 & 65\% & 0.70 \\
\bottomrule
\end{tabular}
\end{table}

The imperfect correlation ($\rho{<}1$) explains why SDGD-cal outperforms SDGD-zero: empirical confusion captures non-semantic factors (prompt position effects, training data frequency) that embeddings miss.


\subsection{SDGD Results on Few-NERD}
\label{sec:sdgd-fewnerd}

Table~\ref{tab:fewnerd-full} presents the full method comparison on Few-NERD (66 types).

\begin{table}[t]
\centering
\small
\caption{Full method comparison on Few-NERD (66 types, Qwen3-32B).}
\label{tab:fewnerd-full}
\begin{tabular}{lccccc}
\toprule
Method & F1 & P & R & API calls & Wall-clock \\
\midrule
Single-pass       & 0.414 & 0.648 & 0.304 & 1  & 2.8 min \\
Random ($k{=}8$)  & 0.436 & 0.530 & 0.370 & 8  & 12.5 min \\
SDGD-zero ($k{=}8$) & 0.469 & 0.548 & 0.410 & 8  & 12.5 min \\
Per-type ($k{=}66$)  & 0.475 & 0.512 & 0.443 & 66 & 68.5 min \\
SDGD-cal ($k{=}9$)  & \textbf{0.488} & 0.556 & 0.435 & 17 & 14.2 min \\
\bottomrule
\end{tabular}
\end{table}

Key observations:
\begin{itemize}
\item \textbf{SDGD-zero vs.\ single-pass}: $+5.5$ F1 pp ($p{=}0.003$), with $+10.6$pp recall improvement offset by $-10.0$pp precision decline. The recall gain reflects reduced attention dilution; the precision decline comes from merge artifacts across 8 batches.
\item \textbf{SDGD-zero vs.\ random batching}: $+3.3$ F1 pp at the same API cost ($k{=}8$). This gap isolates the value of confusability-aware grouping from simple schema decomposition.
\item \textbf{SDGD-cal vs.\ per-type}: $+1.3$ F1 pp with $3.9\times$ fewer API calls. SDGD-cal \emph{exceeds} per-type splitting because it preserves beneficial cross-type context while eliminating high-confusability co-occurrences.
\item \textbf{Type confusion rate}: SDGD-zero reduces type confusion by $44.7\%$ relative to single-pass ($0.021$ vs.\ $0.038$), directly demonstrating Channel 2 mitigation.
\end{itemize}


\subsection{SDGD Results on MAVEN}
\label{sec:sdgd-maven}

Table~\ref{tab:maven-full} shows MAVEN results (168 types).
SDGD uses MiniLM embeddings, which yield significant pairwise signal on MAVEN ($r{=}0.120$), unlike Qwen3 ($r{=}{-}0.007$).

\begin{table}[t]
\centering
\small
\caption{Full comparison on MAVEN (168 types, Qwen3-32B, MiniLM embeddings).}
\label{tab:maven-full}
\begin{tabular}{lccccc}
\toprule
Method & Rel.\ F1 & Exact F1 & API calls & $\Delta$ Rel.\ F1 & Latency \\
\midrule
Single-pass         & 0.272 & 0.014 & 1   & ---     & 1.0$\times$ \\
Random ($k{=}14$)   & 0.282 & 0.017 & 14  & +1.0pp  & 5.3$\times$ \\
SDGD-zero ($k{=}14$)& 0.305 & 0.022 & 14  & +3.3pp  & 5.3$\times$ \\
SDGD-cal            & 0.318 & 0.025 & 24  & +4.6pp  & 8.5$\times$ \\
Per-type ($k{=}168$) & 0.325 & ---   & 168 & +5.3pp  & 40.0$\times$ \\
\bottomrule
\end{tabular}
\end{table}

SDGD-zero significantly outperforms both single-pass ($+3.3$pp, $p{=}0.012$) and random batching ($+2.3$pp, $p{=}0.038$).
SDGD-cal approaches per-type performance ($0.318$ vs.\ $0.325$) with $7\times$ fewer API calls.
Type confusion reduction is $39.7\%$ ($0.035$ vs.\ $0.058$).

The smaller gains on MAVEN vs.\ Few-NERD ($+3.3$ vs.\ $+5.5$pp) are expected: MAVEN's concentrated SCS distribution offers less ``separation space'' for graph coloring, and some count-context benefit is lost when batches shrink from 168 to 12 types.


\subsection{Ablation: Random Batching}
\label{sec:ablation-random}

The random batching ablation (Table~\ref{tab:ablation}) decomposes SDGD's total gain into count-reduction and confusability-aware contributions.

\begin{table}[t]
\centering
\small
\caption{Ablation: SDGD-zero vs.\ random batching vs.\ single-pass on Few-NERD.}
\label{tab:ablation}
\begin{tabular}{lcccc}
\toprule
Method & F1 & Recall & Confusion rate & Mean BCS \\
\midrule
Single-pass ($k{=}1$) & 0.414 & 0.304 & 0.038 & 0.847 \\
Random ($k{=}8$)      & 0.436 & 0.370 & 0.031 & 0.838 \\
SDGD-zero ($k{=}8$)   & 0.469 & 0.410 & 0.021 & 0.762 \\
\bottomrule
\end{tabular}
\end{table}

Of SDGD-zero's total $+5.5$pp F1 gain over single-pass:
\begin{itemize}
\item $+2.2$pp ($40\%$) comes from count reduction alone (random batching vs.\ single-pass).
\item $+3.3$pp ($60\%$) comes from confusability-aware grouping (SDGD-zero vs.\ random).
\end{itemize}

Random batching barely reduces mean batch BCS ($0.838$ vs.\ $0.847$), because random assignment does not systematically separate confusable pairs.
SDGD-zero reduces BCS to $0.762$ ($-10\%$), which translates to a $32\%$ reduction in type confusion rate ($0.021$ vs.\ $0.031$).
SDGD-zero's F1 ($0.469$) exceeds the upper bound of random batching's distribution ($\mu + 2\sigma = 0.436 + 0.024 = 0.460$, over 5 seeds with $\sigma{=}0.012$), confirming the result is not a ``lucky draw.''


\subsection{Ablation: Per-Type Splitting}
\label{sec:ablation-pertype}

Per-type splitting ($k{=}66$) achieves F1 $= 0.475$---only $+0.6$pp above SDGD-zero ($0.469$) but at $8.25\times$ the API cost.
Notably, SDGD-cal ($0.488$) \emph{surpasses} per-type splitting by $+1.3$pp.

Per-type splitting nearly eliminates type confusion (rate $= 0.002$) but introduces new failure modes:
hallucination rate increases to $28\%$ (vs.\ $\sim\!21\%$ in batched conditions) because single-type prompts lack cross-type constraints that suppress spurious predictions.
This reveals that \textbf{optimal decomposition $\neq$ maximum decomposition}: the best strategy balances confusability reduction with context preservation.
SDGD's graph coloring naturally achieves this balance---types are grouped if and only if they are below the confusability threshold.


\subsection{SDGD Variant Comparison}
\label{sec:sdgd-variants}

Table~\ref{tab:sdgd-variants} summarizes SDGD-zero vs.\ SDGD-cal on both datasets.

\begin{table}[t]
\centering
\small
\caption{SDGD variant comparison across datasets.}
\label{tab:sdgd-variants}
\begin{tabular}{llccc}
\toprule
Dataset & Method & F1 & $\Delta$ vs.\ SP & API calls \\
\midrule
\multirow{2}{*}{Few-NERD} & SDGD-zero & 0.469 & +5.5pp & 8 \\
                          & SDGD-cal  & 0.488 & +7.4pp & 17 \\
\midrule
\multirow{2}{*}{MAVEN}    & SDGD-zero & 0.305 & +3.3pp & 14 \\
                          & SDGD-cal  & 0.318 & +4.6pp & 24 \\
\bottomrule
\end{tabular}
\end{table}

SDGD-cal's marginal improvement over SDGD-zero ($+1.9$pp on Few-NERD, $+1.3$pp on MAVEN) comes from replacing the embedding proxy with empirical confusion data.
The Spearman correlation between empirical confusability and Qwen3 SCS is $0.58$ (Few-NERD) and $0.49$ (MAVEN/MiniLM), explaining both why SDGD-zero is effective (the proxy captures most of the signal) and why SDGD-cal improves further (it corrects the $\sim\!15\%$ of misordered pairs).

SDGD-zero captures $74\%$ of SDGD-cal's total gain on Few-NERD ($5.5/7.4$) and $72\%$ on MAVEN ($3.3/4.6$), making it the recommended default for cold-start scenarios without calibration data.


%----------------------------------------------------------------------
\section{Discussion}
\label{sec:discussion}

\paragraph{Practical implications.}
SDGD provides a practical, principled alternative to ad-hoc schema decomposition.
For schemas with ${\geq}30$ types, SDGD-zero offers substantial F1 gains at moderate API cost ($k{=}8$--$14$ calls vs.\ $66$--$168$ for per-type splitting).
When even a handful of calibration examples are available, SDGD-cal further improves performance while remaining efficient.

\paragraph{Embedding selection.}
No single embedding universally captures confusability.
On cross-domain schemas (Few-NERD), instruction-tuned Qwen3 embeddings are superior (P@10 $= 0.90$).
On same-domain homogeneous schemas (MAVEN), generic MiniLM embeddings are more effective.
SDGD-cal sidesteps this choice entirely by using empirical confusion data.
We recommend starting with SDGD-zero using the extraction model's own embeddings, falling back to generic embeddings if the pairwise signal is weak.

\paragraph{Limitations.}
Our pairwise effect sizes are modest ($r{=}0.12$--$0.16$), reflecting the many non-confusability factors that influence extraction quality (prompt design, few-shot examples, input text difficulty).
The $2{\times}2$ interaction effect is only marginally significant ($p{=}0.087$), preventing definitive conclusions about super-additivity.
Our study covers three IE tasks but only two extraction models; further validation on encoder-based LLMs and multilingual settings is warranted.
SDGD's latency overhead ($5$--$8\times$ wall-clock time) may be prohibitive for latency-sensitive applications.

\paragraph{Centrality confound as a methodological contribution.}
Our discovery that per-type confusability aggregation fails on homogeneous schemas due to centrality confounding has implications beyond this work.
Any method that ranks types by average similarity to other types (a common heuristic) risks conflating type centrality with type confusability when the schema lacks clear domain boundaries.
We recommend that future work on schema complexity operate at the pairwise level.


%----------------------------------------------------------------------
\section{Conclusion}
\label{sec:conclusion}

We presented the first controlled study of schema-induced degradation in training-free LLM information extraction, establishing three findings:
(1)~pairwise type confusability independently predicts extraction degradation across NER, event detection, and relation extraction ($r{=}0.12$--$0.16$);
(2)~count and confusability degrade performance through orthogonal channels---omission and type substitution, respectively;
(3)~SDGD, a graph-coloring decomposition framework operating on pairwise confusability, yields $+5.5$ F1 on Few-NERD and $+3.3$ on MAVEN over single-pass extraction, with $60\%$ of the gain attributable to confusability-aware grouping rather than simple count reduction.
Our analysis additionally reveals a centrality confound in per-type confusability aggregation on homogeneous schemas, motivating pairwise approaches.


\bibliographystyle{plainnat}
\begin{thebibliography}{10}

\bibitem[Ding et al.(2021)]{fewnerd}
N.~Ding, G.~Xu, Y.~Chen, et al.
\newblock Few-NERD: A few-shot named entity recognition dataset.
\newblock In \emph{Proc. ACL}, 2021.

\bibitem[Wang et al.(2023)]{maven}
X.~Wang, Z.~Wang, et al.
\newblock MAVEN: A massive general domain event detection dataset.
\newblock In \emph{Proc. EMNLP}, 2020.

\bibitem[Yao et al.(2019)]{docred}
Y.~Yao, D.~Ye, P.~Li, et al.
\newblock DocRED: A large-scale document-level relation extraction dataset.
\newblock In \emph{Proc. ACL}, 2019.

\bibitem[Zhou et al.(2023)]{uniner}
W.~Zhou, S.~Zhang, Y.~Gu, M.~Chen, and H.~Peng.
\newblock UniNER: Universal named entity recognition with instruction-following large language models.
\newblock \emph{arXiv preprint arXiv:2308.03279}, 2023.

\bibitem[Liu et al.(2024)]{lost-middle}
N.~F.~Liu, K.~Lin, J.~Hewitt, A.~Paranjape, M.~Bevilacqua, F.~Petroni, and P.~Liang.
\newblock Lost in the middle: How language models use long contexts.
\newblock \emph{TACL}, 2024.

\bibitem[Wei et al.(2023)]{llmpee}
X.~Wei, J.~Cui, N.~Cheng, et al.
\newblock Zero-shot information extraction via chatting with ChatGPT.
\newblock \emph{arXiv preprint arXiv:2302.10205}, 2023.

\end{thebibliography}

\end{document}
